% =====================================================================
\chapter{Миксеры}
\label{ch:mixes}
% =====================================================================
\chapterintro
  {как входы и переключатели попадают в каналы, как «смешать» два движения в одной серве
   (летающее крыло, V-хвост), как сделать отсечку газа и плавное движение шасси}
  {модель с настроенными входами}

\section{Что делает миксер}

\begin{simple}
Миксер --- это «распределитель»: он решает, \textbf{какой канал} получит какой сигнал.
В простом случае всё один-к-одному: крен → канал 1, тангаж → канал 2. Но иногда одной
серве нужно слушаться \emph{двух} стиков сразу --- тогда миксер складывает сигналы, как
смешивают краски.
\end{simple}

Страница \mpath{MDL\sep Mixes}. Слева --- каналы, справа --- что в них идёт:

\begin{center}
\lcdscreen{MIXES}{7/13}{\lr{CH1 I1 Ail}{100}\lr{CH2 I2 Ele}{100}\lr{CH3 I3 Thr}{100}\lr{CH4 I4 Rud}{100}\lsel{CH5}{}\lblank}
\end{center}

Чтобы добавить строку в пустой канал --- встань на него и нажми \btn{ENTER}. Чтобы добавить
вторую строку в канал --- долгое \btn{ENTER} → \menu{Insert after}.

\section{Переключатели в каналы (AUX)}

Самое частое дело: отправить переключатели на каналы 5, 6, 7…

\begin{figure}[H]
\centering
\begin{tikzpicture}[every node/.style={font=\sffamily\small}]
  \foreach \y/\s/\c/\p/\v/\t in {0/SA/CH5/1/100/арм,-1.3/SB/CH6/0/0/режим полёта,-2.6/SC/CH7/-1/-100/бипер}{
    \pic[transform shape,scale=0.8] at (0,\y) {toggle=\p};
    \node[font=\sffamily\small\bfseries] at (0.9,\y+0.1) {\s};
    \draw[arr] (1.5,\y+0.1) -- (3.3,\y+0.1);
    \node[font=\sffamily\scriptsize] at (2.4,\y+0.35) {миксер};
    \node[draw=etx,fill=etx!8,rounded corners=2pt,minimum width=12mm] at (4.0,\y+0.1) {\c};
    \fill[black!6,draw=black!30] (5.0,\y-0.05) rectangle (8.0,\y+0.25);
    \draw[black!40] (6.5,\y-0.1) -- (6.5,\y+0.3);
    \pgfmathsetmacro{\w}{1.5*\v/100}
    \fill[etx!65] (6.5,\y-0.05) rectangle ++(\w,0.3);
    \node[anchor=west,font=\sffamily\scriptsize] at (8.2,\y+0.1) {\v\quad (\t)};}
\end{tikzpicture}
\caption{Переключатель как источник: вверху $-100$, посередине 0, внизу $+100$}
\end{figure}

\begin{center}\small
\begin{tabular}{@{}llll@{}}
\toprule
Канал & Source & Weight & Зачем \\
\midrule
\sw{CH5} & \sw{SA} & 100 & арм (в ELRS --- только двухпозиционный!) \\
\sw{CH6} & \sw{SB} & 100 & режимы полёта контроллера \\
\sw{CH7} & \sw{SC} & 100 & бипер, «черепаха», запись видео \\
\sw{CH8} & \sw{S1} & 100 & наклон камеры (крутилка) \\
\bottomrule
\end{tabular}
\end{center}

\section{Параметры строки миксера}

\begin{center}
\lcdscreen{MIXER}{CH1}{\lr{Mix name}{}\lr{Source}{I1 Ail}\lsel{Weight}{\ledit{50}}\lr{Offset}{0}\lr{Switch}{---}\lr{Multiplex}{Add}}
\end{center}

\begin{center}\small
\renewcommand{\arraystretch}{1.2}
\begin{tabularx}{\linewidth}{@{}L{30mm}Y@{}}
\toprule
\menu{Source} & что брать: вход (\sw{I1}…), стик, переключатель, \sw{MAX} (всегда $+100$), другой канал \\
\menu{Weight} & вес (минус --- в обратную сторону) \\
\menu{Offset} & сдвиг \\
\menu{Switch} & когда строка работает \\
\menu{Multiplex} & как объединить с верхними строками: \menu{Add} --- \textbf{сложить},
  \menu{Multiply} --- умножить, \menu{Replace} --- \textbf{заменить} \\
\menu{Slow up/down} & плавность: за сколько секунд значение пройдёт весь путь \\
\menu{Delay up/down} & задержка перед началом движения \\
\bottomrule
\end{tabularx}
\end{center}

\section{Летающее крыло: элевоны}

У летающего крыла нет хвоста. Две поверхности сзади --- \textbf{элевоны} --- работают и как
элероны, и как руль высоты:

\begin{figure}[H]
\centering
\begin{tabular}{@{}c@{\hspace{3mm}}c@{\hspace{3mm}}c@{}}
\begin{tikzpicture}[scale=0.62]
  \pic[transform shape,elL/.style={surf},elR/.style={surf}] {wingtop};
  \upmark{(-2.4,-1.35)}\upmark{(2.4,-1.35)}
\end{tikzpicture} &
\begin{tikzpicture}[scale=0.62]
  \pic[transform shape,elL/.style={surf},elR/.style={surf}] {wingtop};
  \downmark{(-2.4,-1.35)}\upmark{(2.4,-1.35)}
\end{tikzpicture} &
\begin{tikzpicture}[scale=0.62]
  \pic[transform shape,elL/.style={surf},elR/.style={surf}] {wingtop};
  \node[circle,fill=black!40,text=white,font=\sffamily\scriptsize\bfseries,inner sep=1.2pt] at (-2.4,-1.35) {0};\upmark{(2.4,-1.35)}
\end{tikzpicture}\\
\small стик на себя: & \small стик вправо: & \small на себя + вправо: \\
\small оба вверх → нос вверх & \small правый вверх, левый вниз & \small правый вверх, левый почти в нуле
\end{tabular}
\caption{Что делают элевоны. ↑ --- поверхность вверх, ↓ --- вниз}
\end{figure}

Значит, каждому элевону нужен и крен, и тангаж --- их надо \textbf{сложить}:

\begin{figure}[H]
\centering
\begin{tikzpicture}[every node/.style={font=\sffamily\small}]
  \node[blk,fill=blkin,minimum width=20mm] (ail) at (0,1) {\sw{I1} крен};
  \node[blk,fill=blkin,minimum width=20mm] (ele) at (0,-1) {\sw{I2} тангаж};
  \node[circle,draw=etx,fill=blkmix,minimum size=9mm,font=\sffamily\Large\bfseries] (p1) at (4,1) {+};
  \node[circle,draw=etx,fill=blkmix,minimum size=9mm,font=\sffamily\Large\bfseries] (p2) at (4,-1) {+};
  \node[blk,fill=blkout,minimum width=26mm] (c1) at (7.3,1) {\sw{CH1} левый элевон};
  \node[blk,fill=blkout,minimum width=26mm] (c2) at (7.3,-1) {\sw{CH2} правый элевон};
  \draw[arr] (ail) -- node[above,lbl]{$\times(+50\,\%)$} (p1);
  \draw[arr] (ele) -- node[below right,lbl,pos=0.82]{$\times(+50\,\%)$} (p1);
  \draw[arr] (ail) -- node[above right,lbl,pos=0.82]{$\times(-50\,\%)$} (p2);
  \draw[arr] (ele) -- node[below,lbl]{$\times(+50\,\%)$} (p2);
  \draw[arr] (p1) -- (c1);
  \draw[arr] (p2) -- (c2);
\end{tikzpicture}
\caption{Каждый элевон = половина крена + половина тангажа. У правого крен со знаком минус}
\end{figure}

\begin{recipe}{элевоны}
\begin{center}\small
\begin{tabular}{@{}llll@{}}
\toprule
Канал & Source & Weight & Multiplex \\
\midrule
\sw{CH1} & \sw{I1} Ail & $+50$ & Add \\
         & \sw{I2} Ele & $+50$ & Add \\
\sw{CH2} & \sw{I1} Ail & $-50$ & Add \\
         & \sw{I2} Ele & $+50$ & Add \\
\sw{CH3} & \sw{I3} Thr & 100 & Add \\
\bottomrule
\end{tabular}
\end{center}
\end{recipe}

\begin{whyit}
Посчитаем. Стик крена вправо на 40\,\% (это $+40$), стик тангажа на себя на 60\,\%
(на себя --- это минус: $-60$).
\begin{itemize}
  \item \sw{CH1} (левый) $= 40\cdot0{,}5 + (-60)\cdot0{,}5 = 20 - 30 = -10$;
  \item \sw{CH2} (правый) $= 40\cdot(-0{,}5) + (-60)\cdot0{,}5 = -20 - 30 = -50$.
\end{itemize}
Пусть минус у нас означает «поверхность вверх». Тогда оба элевона идут вверх --- нос
поднимается, но правый сильно, а левый чуть-чуть --- крыло кренится вправо. Всё как надо!
(Если у твоей сервы «вверх» --- это плюс, это исправляется реверсом на странице \menu{Outputs}.)
Веса по 50\,\% нужны, чтобы сумма никогда не превысила 100\,\%.
\end{whyit}

Проверяй на земле: стик на себя --- оба элевона вверх; стик вправо --- правый вверх, левый вниз.
Если какой-то ходит наоборот --- включи реверс этого канала на странице \menu{Outputs}
(глава~\ref{ch:outputs}).

\section{V-образный хвост}

Тот же приём, только смешиваются тангаж и рыскание:

\begin{figure}[H]
\centering
\begin{tabular}{@{}c@{\hspace{12mm}}c@{}}
\begin{tikzpicture}[scale=1.1]\pic[transform shape] {vtail={45}{135}};\end{tikzpicture} &
\begin{tikzpicture}[scale=1.1]\pic[transform shape] {vtail={45}{-45}};\end{tikzpicture}\\
\small тангаж «нос вверх»: обе вверх & \small рыскание вправо: обе вправо
\end{tabular}
\caption{V-хвост, вид сзади}
\end{figure}

\sw{CH2} = Ele $+50$ и Rud $+50$; \sw{CH4} = Ele $+50$ и Rud $-50$.

\section{Отсечка газа: строка «Replace»}

\begin{simple}
Строки канала работают по очереди сверху вниз. Строка с \menu{Replace} говорит: «забудь всё,
что насчитали выше, теперь значение --- вот такое». Если её включать переключателем, получится
надёжная «кнопка стоп» для мотора.
\end{simple}

\begin{figure}[H]
\centering
\begin{tikzpicture}[every node/.style={font=\sffamily\small}]
  \node[blk,fill=blkin,minimum width=30mm] (a) at (0,0.6) {строка 1: \sw{I3} газ};
  \node[blk,fill=warn!15,draw=warn,minimum width=30mm] (b) at (0,-0.6) {строка 2: \sw{MAX} $-100$\\\scriptsize Replace, если \sw{SD↓}};
  \node[blk,fill=blkout,minimum width=18mm] (c) at (5.5,0) {\sw{CH3}};
  \draw[arr] (a.east) -- (c.west);
  \draw[arr,warn] (b.east) -- (c.west);
  \node[anchor=west,font=\sffamily\scriptsize,align=left] at (6.5,0.4) {\sw{SD↑}: газ от стика};
  \node[anchor=west,font=\sffamily\scriptsize,align=left,text=warn] at (6.5,-0.4) {\sw{SD↓}: всегда $-100$ (мотор стоп)};
\end{tikzpicture}
\caption{Отсечка газа}
\end{figure}

\begin{recipe}{отсечка газа для самолёта}
В канале \sw{CH3} после строки газа добавь строку: Source = \sw{MAX}, Weight = $-100$,
Switch = \sw{SD↓}, Multiplex = \menu{Replace}. То же можно сделать специальной функцией
\menu{Override} (глава~\ref{ch:special}).
\end{recipe}

\section{Плавность (Slow) и задержка (Delay)}

\begin{figure}[H]
\centering
\begin{tikzpicture}
\begin{axis}[width=11cm,height=4.6cm,xmin=0,xmax=6,ymin=-110,ymax=110,
  xlabel={секунды},ylabel={канал},grid=major,grid style={gray!20},
  tick label style={font=\scriptsize},label style={font=\small},legend pos=south east,legend style={font=\scriptsize}]
\addplot[very thick,black!50,const plot] coordinates {(0,-100) (1,100) (6,100)};
\addplot[very thick,accent] coordinates {(0,-100) (1,-100) (3,100) (6,100)};
\addplot[very thick,etx,dashed] coordinates {(0,-100) (2,-100) (4,100) (6,100)};
\legend{переключатель,Slow 2\,с,Delay 1\,с + Slow 2\,с}
\end{axis}
\end{tikzpicture}
\caption{Щёлкнул переключатель в момент 1\,с. С \menu{Slow} канал плавно доедет за 2 секунды,
с \menu{Delay} --- ещё и начнёт позже}
\end{figure}

Примеры: шасси и закрылки выезжают плавно, «как у настоящих» (\menu{Slow} 2--3\,с);
у планера мотор запускается плавно (\menu{Slow up} 1\,с на газу).

\section{Компенсация при выпуске закрылков}

При выпуске закрылков самолёт задирает нос. Добавь в канал руля высоты строку:
Source = переключатель закрылков, Weight = $-10$ (подбирается в полёте),
Curve = \menu{x>0} --- поправка будет работать только при выпущенных закрылках.

\begin{tryit}
\begin{enumerate}
  \item Отправь \sw{SA} в \sw{CH5}, \sw{SB} в \sw{CH6}. Проверь на странице каналов.
  \item Сделай модель «летающее крыло» с элевонами. Подвигай стиками и посмотри, как меняются
        \sw{CH1} и \sw{CH2}. Посчитай в уме, какие должны быть значения.
  \item Добавь \menu{Slow} 3\,с на \sw{CH6} и посмотри, как медленно ползёт полоска.
\end{enumerate}
\end{tryit}

\begin{summary}
\begin{itemize}
  \item Миксер решает, какой канал получает какой сигнал.
  \item Строки одного канала выполняются сверху вниз: \menu{Add} --- сложить,
        \menu{Replace} --- заменить.
  \item Элевоны и V-хвост --- по две строки с весами 50\,\% и разными знаками.
  \item \menu{Slow} --- плавность, \menu{Delay} --- задержка.
\end{itemize}
\end{summary}
